\documentclass[11pt,a4paper]{article}
\usepackage{fontspec,amsmath,amssymb,amsthm,microtype}
\setmainfont{lmroman10-regular.otf}[BoldFont=lmroman10-bold.otf,ItalicFont=lmroman10-italic.otf,BoldItalicFont=lmroman10-bolditalic.otf]
\usepackage{xcolor}
\usepackage[most]{tcolorbox}
\usepackage{mathrsfs}
\providecommand{\lt}{<}
\providecommand{\gt}{>}
\setmonofont{DejaVuSansMono.ttf}[BoldFont=DejaVuSansMono-Bold.ttf,ItalicFont=DejaVuSansMono-Oblique.ttf,BoldItalicFont=DejaVuSansMono-BoldOblique.ttf,Scale=0.86]
\newfontfamily\leanfallback{STIXTwoMath-Regular.otf}[Scale=0.86]
\definecolor{leanrule}{gray}{0.62}
\definecolor{leanback}{gray}{0.962}
\definecolor{leantitle}{gray}{0.22}
\definecolor{leanmuted}{gray}{0.38}
\newtcolorbox{leanbox}[1]{enhanced,breakable,colback=leanback,colframe=leanrule,
  boxrule=0.4pt,arc=0.6mm,left=2.4mm,right=2.4mm,top=1.8mm,bottom=1.8mm,
  before skip=0.8em,after skip=1em,fonttitle=\small\bfseries,coltitle=leantitle,
  colbacktitle=leanback,titlerule=0.3pt,title={#1}}
\newlength{\leancharwidth}
\newenvironment{leancode}{\par\smallskip\begingroup\ttfamily\footnotesize\raggedright
  \settowidth{\leancharwidth}{M}\setlength{\parskip}{0pt}\setlength{\parindent}{0pt}
  \frenchspacing}{\par\endgroup\smallskip}
\newcommand{\leanline}[2]{\hangindent=\dimexpr#1\leancharwidth+4\leancharwidth\relax
  \hangafter=1\noindent\hspace*{#1\leancharwidth}#2\par}
\newcommand{\leansource}[1]{\par{\footnotesize\color{leanmuted}#1}\par}
\newcommand{\leanlabel}[1]{\par\smallskip\noindent{\small\bfseries #1}\par}
\usepackage[margin=25mm]{geometry}
\usepackage[unicode,hidelinks]{hyperref}
\hypersetup{pdfauthor={LangevinLean Project}}
\setlength{\emergencystretch}{3em}
\setlength{\parskip}{0.4em}
\setlength{\parindent}{0pt}
\allowdisplaybreaks[2]
\raggedbottom
\title{\texorpdfstring{Conditional averaging over an independent path input}{Conditional averaging over an independent path input}}
\author{LangevinLean Project}
\date{}
\begin{document}\maketitle
Draft mathematical exposition; these notes may contain mistakes.


Let \(X\) be measurable from a past sigma-algebra \(\mathcal F\), and let \(Z\) be independent of all of \(\mathcal F\). For an integrable measurable vector function \(F\), conditioning \(F(X,Z)\) on the past integrates only the independent input. The input space may be a whole path space.


\section{\texorpdfstring{Factorization after restriction to a past event}{Factorization after restriction to a past event}}


\textbf{Claim.} For every \(A\in\mathcal F\), the joint law under the restricted measure satisfies

\[\operatorname{Law}_{Q|_A}(X,Z)=\operatorname{Law}_{Q|_A}(X)\otimes\operatorname{Law}_{Q}(Z).\]

\textbf{Proof.} Evaluate both sides on a measurable rectangle. Its preimage is the intersection of a past event, obtained from \(X\) and \(A\), with an event determined by \(Z\). Independence factorizes its mass. Equality on all rectangles determines the product measure.

\begin{leanbox}{Lean: \texttt{independent\_\allowbreak{}input\_\allowbreak{}restrict\_\allowbreak{}joint\_\allowbreak{}map}}
\leansource{\texttt{independent\_\allowbreak{}input\_\allowbreak{}restrict\_\allowbreak{}joint\_\allowbreak{}map} \quad WassersteinMixing/\allowbreak{}ConditionalIndependentAverage.lean, lines 19–35}
\begin{leancode}
\leanline{0}{theorem\ independent\_input\_restrict\_joint\_map\ \{Ω\ A\ B\ :\ Type*\}}
\leanline{4}{(m\ :\ MeasurableSpace\ Ω)\ [m0\ :\ MeasurableSpace\ Ω]\ [MeasurableSpace\ A]\ [MeasurableSpace\ B]}
\leanline{4}{(Q\ :\ Measure\ Ω)\ [IsProbabilityMeasure\ Q]\ (hm\ :\ m\ ≤\ m0)}
\leanline{4}{(X\ :\ Ω\ →\ A)\ (Y\ :\ Ω\ →\ B)\ (hX\ :\ Measurable[m]\ X)\ (hY\ :\ Measurable\ Y)}
\leanline{4}{(hi\ :\ Indep\ ((inferInstance\ :\ MeasurableSpace\ B).comap\ Y)\ m\ Q)}
\leanline{4}{(s\ :\ Set\ Ω)\ (hs\ :\ MeasurableSet[m]\ s)\ :}
\leanline{4}{(Q.restrict\ s).map\ (fun\ ω\ ↦\ (X\ ω,\ Y\ ω))\ =}
\leanline{6}{((Q.restrict\ s).map\ X).prod\ (Q.map\ Y)}
\end{leancode}
\leanlabel{How the Lean proof follows the argument}
\texttt{independent\_\allowbreak{}input\_\allowbreak{}restrict\_\allowbreak{}joint\_\allowbreak{}map} uses product-measure extensionality and the definition of independence on the two preimage events. The independent input retains its unrestricted marginal.
\end{leanbox}

\section{\texorpdfstring{The genuine conditional mean depends only on the current state}{The genuine conditional mean depends only on the current state}}


\textbf{Claim.}

\[\mathbb E[F(X,Z)\mid\mathcal F]=\int F(X,z)\,\operatorname{Law}(Z)(dz).\]

\textbf{Proof.} Product integrability gives integrability of the averaged function and its measurable dependence on \(X\). On every past event, the restricted pair-law factorization and Fubini identify its integral with the integral of \(F(X,Z)\). These are the defining properties of conditional expectation. If \(X\) is only almost surely past-measurable, replace it by a measurable version, apply the result and transport back through almost-sure equality. The conditional mean is then measurable from \(X\) alone.

\begin{leanbox}{Lean: \texttt{condExp\_\allowbreak{}independent\_\allowbreak{}input\_\allowbreak{}average}, \texttt{condExp\_\allowbreak{}independent\_\allowbreak{}input\_\allowbreak{}state\_\allowbreak{}measurable}}
\leansource{\texttt{condExp\_\allowbreak{}independent\_\allowbreak{}input\_\allowbreak{}average} \quad WassersteinMixing/\allowbreak{}ConditionalIndependentAverage.lean, lines 39–80}
\begin{leancode}
\leanline{0}{theorem\ condExp\_independent\_input\_average\ \{Ω\ A\ B\ H\ :\ Type*\}}
\leanline{4}{(m\ :\ MeasurableSpace\ Ω)\ [m0\ :\ MeasurableSpace\ Ω]\ [MeasurableSpace\ A]\ [MeasurableSpace\ B]}
\leanline{4}{[NormedAddCommGroup\ H]\ [NormedSpace\ ℝ\ H]\ [CompleteSpace\ H]}
\leanline{4}{(Q\ :\ Measure\ Ω)\ [IsProbabilityMeasure\ Q]\ (hm\ :\ m\ ≤\ m0)}
\leanline{4}{(X\ :\ Ω\ →\ A)\ (Y\ :\ Ω\ →\ B)\ (hX\ :\ Measurable[m]\ X)\ (hY\ :\ Measurable\ Y)}
\leanline{4}{(hi\ :\ Indep\ ((inferInstance\ :\ MeasurableSpace\ B).comap\ Y)\ m\ Q)}
\leanline{4}{(F\ :\ A\ ×\ B\ →\ H)\ (hF\ :\ StronglyMeasurable\ F)}
\leanline{4}{(hFi\ :\ Integrable\ (fun\ ω\ ↦\ F\ (X\ ω,\ Y\ ω))\ Q)\ :}
\leanline{4}{Q[(fun\ ω\ ↦\ F\ (X\ ω,\ Y\ ω))\ |\ m]\ =ᵐ[Q]}
\leanline{6}{fun\ ω\ ↦\ ∫\ y,\ F\ (X\ ω,\ y)\ ∂Q.map\ Y}
\end{leancode}
\leansource{\texttt{condExp\_\allowbreak{}independent\_\allowbreak{}input\_\allowbreak{}state\_\allowbreak{}measurable} \quad WassersteinMixing/\allowbreak{}ConditionalIndependentAverage.lean, lines 85–112}
\begin{leancode}
\leanline{0}{theorem\ condExp\_independent\_input\_state\_measurable\ \{Ω\ A\ B\ H\ :\ Type*\}}
\leanline{4}{(m\ :\ MeasurableSpace\ Ω)\ [m0\ :\ MeasurableSpace\ Ω]}
\leanline{4}{[NormedAddCommGroup\ A]\ [mA\ :\ MeasurableSpace\ A]\ [BorelSpace\ A]}
\leanline{4}{[MeasurableSpace\ B]\ [NormedAddCommGroup\ H]\ [NormedSpace\ ℝ\ H]\ [CompleteSpace\ H]}
\leanline{4}{(Q\ :\ Measure\ Ω)\ [IsProbabilityMeasure\ Q]\ (hm\ :\ m\ ≤\ m0)}
\leanline{4}{(X\ :\ Ω\ →\ A)\ (Y\ :\ Ω\ →\ B)\ (hX\ :\ AEStronglyMeasurable[m]\ X\ Q)\ (hY\ :\ Measurable\ Y)}
\leanline{4}{(hi\ :\ Indep\ ((inferInstance\ :\ MeasurableSpace\ B).comap\ Y)\ m\ Q)}
\leanline{4}{(F\ :\ A\ ×\ B\ →\ H)\ (hF\ :\ StronglyMeasurable\ F)}
\leanline{4}{(hFi\ :\ Integrable\ (fun\ ω\ ↦\ F\ (X\ ω,\ Y\ ω))\ Q)\ :}
\leanline{4}{AEStronglyMeasurable[mA.comap\ X]\ (Q[(fun\ ω\ ↦\ F\ (X\ ω,\ Y\ ω))\ |\ m])\ Q}
\end{leancode}
\leanlabel{How the Lean proof follows the argument}
\texttt{cond\allowbreak{}Exp\_\allowbreak{}independent\_\allowbreak{}input\_\allowbreak{}average} applies the set-integral characterization of conditional expectation. The final theorem uses the measurable version furnished by \texttt{AEStrongly\allowbreak{}Measurable} and proves measurability for the sigma-algebra generated by the original state.
\end{leanbox}

\end{document}
